\subsection{Limits of norm compatibility and tower descent}
\label{sec:norm-compatibility-towers}

\paragraph{Finding A51 (compatible classes and rational generator bounds).}
Let $C/k$ be smooth, projective and geometrically connected,
of genus $g\ge2$, where $k=\mathbf F_q$ has characteristic $p$.
Suppose hypothetically that distinct $P,Q\in C(k)$ have equal
geometrically pro-prime-to-$p$ point sections. The covering
implication and independent uniqueness argument of recovered
A1, A.31, give compatible rational lifts on every permitted
pointed cover and classes
\[
       \delta_D=[Q_D-P_D]\in J_D(k)[p^\infty]
       \tag{NC.1}
\]
compatible under norms \cite[claim A.31]{A1run}.
The refinement and point-count conclusions below concern these
actual classes. Separately, a Frobenius-stable rank-one
$\mathbf Z_\ell$ quotient, when it exists, gives the stated
uniform bound on the number of rational $p$-primary generators.
The abstract examples prove limits of those algebraic
implications; none is asserted to be a geometric tower or
a point-section collision. No universal witness is supplied.

\begin{proof}
\emph{1. Exact refinement and order formulas.}
Let $h:D'\to D$ be a finite etale refinement of degree $d$
prime to $p$, with the compatible rational point lifts.
Write $\delta'=\delta_{D'}$ and $\delta=\delta_D$.
Pushforward of the point divisors and pullback of a divisor
give
\[
                     h_*\delta'=\delta,
                     \qquad h_*h^*=[d].
       \tag{NC.2}
\]
The second identity follows on each closed point from the
degree-$d$ fibre formula, and hence on divisor classes.
Since $[d]$ is invertible on $p$-primary groups, the map
$d^{-1}h^*$ is a section of $h_*$ on these groups. Therefore
\[
\begin{aligned}
 J_{D'}(k)[p^\infty]
   &=d^{-1}h^*J_D(k)[p^\infty]\oplus K_h,\\
 K_h&=\ker\bigl(h_*:J_{D'}(k)[p^\infty]
                         \longrightarrow J_D(k)[p^\infty]\bigr),\\
 \delta'&=d^{-1}h^*\delta+b,\qquad b\in K_h.
\end{aligned}
       \tag{NC.3}
\]
Indeed subtracting $d^{-1}h^*h_*$ from the identity gives
the second summand, and applying $h_*$ proves that the
intersection is zero. If $m,m'$ are the exact orders of
$\delta,\delta'$, respectively, this direct sum proves
\[
           m\mid m',\qquad
           m'=\operatorname{lcm}\bigl(m,\operatorname{ord}(b)\bigr).
       \tag{NC.4}
\]
These formulas preserve the old order and leave the new
norm-kernel component. They do not assert that every possible
$b$ is represented by compatible point lifts.

\emph{2. The principal-function norm relation.}
A geometric function with divisor $m(Q_D-P_D)$ descends,
after multiplication by a scalar, to $f\in k(D)^\times$.
Its Galois conjugates differ by constants, and Hilbert90
removes that scalar cocycle
\cite[Theorem~59.24.1, Tag~\texttt{03P8}]{Stacks03P8}.
Likewise choose $f'\in k(D')^\times$ with divisor
$m'(Q_{D'}-P_{D'})$.
The functions are separable: if one were a $p$th power over
the perfect field, its divisor would contradict the exact
order; for order one, its simple poles already exclude that
possibility. The divisor-of-norm formula gives
\[
                 N_h(f')=c\,f^{m'/m},\qquad c\in k^\times
       \tag{NC.5}
\]
\cite[Lemma~42.18.1, Tag~\texttt{02RT}]{Stacks02RT}.
Here the rational points have residue degree one, so their
pushforwards have precisely the displayed multiplicities.
The quotient of the two sides has zero divisor on the
proper geometrically connected curve, hence is constant.

A $p$th-power norm does not, by itself, force a $p$th-power
function upstairs. To check this directly, choose $d>1$
prime to $p$ and a finite field containing a nontrivial
$d$th root $\zeta$. In $k(z)/k(t)$ with $t=z^d$,
\[
                    N\!\left(\frac{z-1}{z-\zeta}\right)=1.
       \tag{NC.6}
\]
The function has a simple zero and pole and is not a $p$th
power. This is only an algebraic counterexample to the
unrestricted norm inference: the map is ramified, and its
zero and pole have the same image. It is not an example of
the pointed etale-cover configuration in (NC.1).

\emph{3. The split-point estimate is a lower bound.}
Suppose $D\to C$ has degree $n$, is geometrically Galois,
and has constant deck group over $k$. A rational lift of
$P$ makes its whole fibre rational; the lift of $Q$ does
the same for its disjoint fibre. Thus $\#D(k)\ge2n$.
The function $f$ of Step2 has degree $m$. Its zero and pole
fibres contain one point each, and each other rational value
has at most $m$ rational preimages. Consequently
\[
     2n\le\#D(k)\le(q-1)m+2,
     \qquad
     m\ge\left\lceil\frac{2n-2}{q-1}\right\rceil.
       \tag{NC.7}
\]
This estimate concerns the particular point-difference order.
The constant-deck hypothesis is essential for the $2n$
rational points. No existence of arbitrarily large constant
pointed covers over this fixed field is asserted, and
enlarging the field changes the bound.

\emph{4. Proper-quotient norms miss the faithful component.}
Let $D_i\to C$ be constant cyclic covers of coprime odd
degrees $n_i>1$ prime to $p$, both with rational lifts of $P$.
Their fibre product $D$ is geometrically connected: the two
geometric Galois fields have intersection of degree dividing
both $n_1,n_2$. Its group is
$G=C_{n_1}\times C_{n_2}\simeq C_n$, $n=n_1n_2$.
For $h_i:D\to D_i$, set $K_i=\ker(G\to C_{n_i})$.
The pullback of the pushed-forward divisor is its orbit sum,
so
\[
 h_i^*h_{i*}=\sum_{\gamma\in K_i}\gamma,
 \qquad
 \left(\sum_{\gamma\in K_i}\gamma\right)p_n=0,
 \qquad p_n=ne_n.
       \tag{NC.8}
\]
The second identity holds in $\mathbf Z[G]$: a faithful
character is nontrivial on $K_i$, so its character sum is
zero, while $p_n$ kills all nonfaithful characters.
Character evaluation is faithful on $\mathbf Q[G]$.
For the actual primitive subvariety $P_n=\operatorname{im}p_n$
of A46, $p_n$ acts as $[n]$. Since $n$ and $\deg h_i$ are
prime to $p$, (NC.2) and (NC.8) show
\[
                       h_{i*}\bigl(P_n(k)[p^\infty]\bigr)=0.
       \tag{NC.9}
\]
Thus both old norm constraints are insensitive to the full
faithful product-conductor component. The same proof applies
to every proper cyclic quotient of a fixed cyclic cover.
A good packet downstairs constrains its inflated characters,
not the faithful characters killed by that norm. This is a
statement about actual Jacobian maps. It does not assert that
an arbitrary undetected torsion point can still be represented
by the compatible two-point divisor, or that no different
argument can combine the packet information.

\emph{5. The minimal deck field in a rank-one tower.}
Fix an odd prime $\ell\ne p$. Suppose a Frobenius-stable
rank-one $\mathbf Z_\ell$ quotient of
$\pi_1(C_{\bar k})^{\rm ab}$ is given. Let its Frobenius
multiplier be $a\in\mathbf Z_\ell^\times$, using arithmetic
Frobenius on geometric points. The pointed section gives
the associated compatible geometrically cyclic covers
$D_r\to C$ of degree $\ell^r$ over $k$.
The Abel--Jacobi identification with the Tate module
\cite[Proposition~9.1]{MilneJV}
\cite[Remark~15.3]{MilneAV}
shows that $a$ is a Frobenius eigenvalue on $T_\ell J_C$.
The Weil bound on its complex absolute values makes it
a non-root of unity \cite[Theorem~19.1]{MilneAV}.
Set
\[
 t=\operatorname{ord}_{\ell}(a),\qquad
 s=v_\ell(a^t-1),\qquad
 f_r=\operatorname{ord}_{\ell^r}(a).
       \tag{NC.10}
\]
Then $t\mid\ell-1$, $1\le s<\infty$, and
\[
                         f_r=t\ell^{\max(0,r-s)}.
       \tag{NC.11}
\]
To prove this formula, write $a^t=1+\ell^s u$ with $u$ a
unit. A power prime to $\ell$ preserves the valuation $s$;
an $\ell$th power raises it by exactly one. Both assertions
follow from binomial expansion, since $\ell$ is odd and
the first nonconstant term has strictly smaller valuation
than the remaining terms. Iteration gives
$v_\ell((a^t)^h-1)=s+v_\ell(h)$.
Any exponent yielding one modulo $\ell$ is a multiple of
$t$, proving (NC.11). This is exactly the minimal extension
degree defining all geometric deck transformations of $D_r$.
The assertion is conditional on the stable rank-one quotient;
it does not claim that every chosen $\ell$ supplies one.

\emph{6. A uniform rational generator bound, with no exponent bound.}
Let $B_r$ be the actual primitive conductor-$\ell^r$
subvariety of $J_{D_r}$ and let $k_j=\mathbf F_{q^j}$.
Use the full finite etale algebra, with its indicated action,
\[
\begin{aligned}
 \mathcal O_r&=\mathbf Z_p[Z]/\Phi_{\ell^r}(Z),\\
 \tau_r(Z)&=Z^a,\qquad R_r=\mathcal O_r^{\tau_r},\\
 U_r&=B_r(k_{f_r})[p^\infty],\qquad
 N_r=B_r(k)[p^\infty].
\end{aligned}
       \tag{NC.12}
\]
The integral descent result (SD.3)--(SD.10) in A45 applies
to $U_r$: its Frobenius is $\tau_r$-semilinear and has
$f_r$th power one. It gives
\[
                 U_r\simeq\mathcal O_r\otimes_{R_r}N_r.
       \tag{NC.13}
\]
Its known descent input is effective faithfully flat descent
for arbitrary modules, not averaging by $1/f_r$
\cite[Proposition~35.3.9, Tag~\texttt{023N}]{Stacks023N}.
The finite etale $R_r$-algebra $\mathcal O_r$ has constant
rank $f_r$. Since $R_r$ is a product of unramified discrete
valuation rings, this module is free of rank $f_r$, component
by component. Hence (NC.13) is, additively and noncanonically,
a direct sum of $f_r$ copies of $N_r$, and
\[
\begin{aligned}
 \#U_r&=(\#N_r)^{f_r},\qquad \exp U_r=\exp N_r,\\
 \dim_{\mathbf F_p}U_r[p]&=f_r\dim_{\mathbf F_p}N_r[p].
\end{aligned}
       \tag{NC.14}
\]
The dimension of $B_r$ is $(g-1)\varphi(\ell^r)$: for a
prime-power cover it is the new factor between $D_r$ and
$D_{r-1}$, so this follows by subtracting their etale
Riemann--Hurwitz genus formulas. The geometric $p$-rank
of an abelian variety is at most its dimension
\cite[Remark~8.5, the bound $r\le g$]{MilneAV}.
Thus (NC.14) gives
\[
\begin{aligned}
 \dim_{\mathbf F_p}B_r(k)[p]
 &\le \frac{(g-1)\varphi(\ell^r)}{f_r}\\
 &=\frac{(g-1)(\ell-1)\ell^{s-1}}{t}
 \qquad(r\ge s).
\end{aligned}
       \tag{NC.15}
\]
This is a uniform bound on the number of rational
$p$-primary generators.

For every $e\ge1$, however, the \emph{abstract, nongeometric}
module $\mathcal O_r/p^e\mathcal O_r$ with Frobenius
$\tau_r$ satisfies exactly the same semilinear descent
hypotheses. Its invariants are $R_r/p^eR_r$, by that descent,
and have exponent $p^e$. Thus the displayed descent identities
and generator bound do not by themselves supply an exponent
bound or vanishing. No such module is asserted to be a
Jacobian factor or part of a geometric tower.

\emph{7. Finite invariant groups also allow arbitrary abstract exponents.}
There is a compatible \emph{abstract, nongeometric} example
with one fixed group action. Let $G=\mathbf Z_\ell$, choose
arbitrary integers $e_r\ge1$, and put
\[
 A=\bigoplus_{r\ge1}\mathcal O_r/p^{e_r}\mathcal O_r,
 \qquad H_m=\ell^m\mathbf Z_\ell.
       \tag{NC.16}
\]
Here a topological generator of $G$ acts on the $r$th
summand by multiplication by $Z$, not by $\tau_r$.
Every element has finite support and an open stabilizer,
so this action on the discrete group is continuous.
For $r\le m$ the summand is fixed by $H_m$.
For $r>m$, the element $Z^{\ell^m}-1$ is a unit in
$\mathcal O_r$: its value at every primitive $\ell^r$th
root in characteristic $p$ is nonzero, since $p\ne\ell$.
It is also a unit modulo $p^{e_r}$, so that summand has no
invariants. Consequently
\[
\begin{aligned}
 A^{H_m}&=\bigoplus_{1\le r\le m}
                 \mathcal O_r/p^{e_r}\mathcal O_r,\\
 \exp(A^{H_m})&=p^{\max_{1\le r\le m}e_r}\qquad(m\ge1).
\end{aligned}
       \tag{NC.17}
\]
Every open-subgroup invariant group is finite, while these
exponents can grow arbitrarily. This is solely an algebraic
limitation; no curve, tower of Jacobians or compatible point
divisors realizing $A$ are claimed.
\end{proof}

The unresolved covering-route input is still an upper bound
on a suitable projected \emph{point-difference} order, or a
geometric argument forcing its faithful component to vanish
from the signed-orbit constraints. Norm compatibility supplies
neither conclusion through the formulas proved here. This
records exact consequences and limits, without turning the
absence of a theorem in a bounded source search into a theorem
of nonexistence, and without a claim of main-gate reduction.
